\chapter{Başlangıç ve Son Topoloji}
\label{ch:initial_final}

\textbf{Başlangıç topolojisi} (initial topology), verilen haritaları sürekli kılan en
kaba (coarsest) topolojidir. \textbf{Son topoloji} (final topology), verilen haritaları
sürekli kılan en ince (finest) topolojidir. Altuzay ve çarpım topolojileri başlangıç
topolojisinin; bölüm topolojisi ise son topolojisinin özel halleridir.

\begin{sezgi}
Başlangıç topolojisini, $X$'e ``tam yeterince'' açık küme koyan bir ayar düğmesi
gibi düşünün. Çok az açık koyarsanız $f_\alpha$ haritaları sürekli olmaz; çok fazla
koyarsanız gereksizdir. Geri çekme $f_\alpha^{-1}(U)$ ile üretilen alt-baz,
haritaları sürekli kılan \emph{en küçük} (en kaba) açık-küme koleksiyonudur. Son
topoloji ise tam tersi yönde çalışır: $Y$'ye en çok açık kümeyi koyar, ama her
$g_i$'yi sürekli bozmadan.
\end{sezgi}

\begin{figure}[ht]
\centering
\begin{tikzpicture}[font=\small]
  % Domain X (left)
  \draw[rounded corners=6pt, thick] (-0.6,-1.5) rectangle (1.9,1.6);
  \node at (0.65,1.85) {$X$ (taşıyıcı, topolojisiz)};
  \draw[blue!70!black, thick, fill=blue!10]
    (0.65,0.55) ellipse [x radius=0.78, y radius=0.42];
  \node[blue!70!black, font=\scriptsize] at (0.65,0.55) {$f_1^{-1}(U_1)$};
  \draw[red!70!black, thick, fill=red!10]
    (0.65,-0.7) ellipse [x radius=0.78, y radius=0.42];
  \node[red!70!black, font=\scriptsize] at (0.65,-0.7) {$f_2^{-1}(U_2)$};

  % Codomain Y_1 (top right)
  \draw[rounded corners=6pt, thick] (5.4,0.5) rectangle (7.6,1.9);
  \node at (6.5,2.15) {$(Y_1,\tau_1)$};
  \draw[blue!70!black, thick, fill=blue!12]
    (6.5,1.2) ellipse [x radius=0.55, y radius=0.32];
  \node[blue!70!black, font=\scriptsize] at (6.5,1.2) {$U_1$};

  % Codomain Y_2 (bottom right)
  \draw[rounded corners=6pt, thick] (5.4,-1.9) rectangle (7.6,-0.5);
  \node at (6.5,-2.2) {$(Y_2,\tau_2)$};
  \draw[red!70!black, thick, fill=red!12]
    (6.5,-1.2) ellipse [x radius=0.55, y radius=0.32];
  \node[red!70!black, font=\scriptsize] at (6.5,-1.2) {$U_2$};

  % Maps
  \draw[blue!70!black, -{Stealth[length=2.4mm]}, thick]
    (1.5,0.55) to[bend left=12] node[above, font=\scriptsize, pos=0.55] {$f_1$} (5.95,1.2);
  \draw[red!70!black, -{Stealth[length=2.4mm]}, thick]
    (1.5,-0.7) to[bend right=12] node[below, font=\scriptsize, pos=0.55] {$f_2$} (5.95,-1.2);

  \node[font=\scriptsize, align=center] at (3.0,-3.0)
    {Alt-baz $\mathcal{S}=\{f_\alpha^{-1}(U)\}$ geri çekme ile üretilir;\\
     $\tau_{\mathrm{ini}}=\langle\mathcal{S}\rangle$ her $f_\alpha$'yı sürekli kılan \emph{en kaba} topoloji.};
\end{tikzpicture}

\caption{Başlangıç topolojisi: bir harita ailesini sürekli kılan en kaba topoloji,
geri çekme $f_\alpha^{-1}(U)$ ile üretilen alt-bazdan doğar.}
\label{fig:ch13_baslangic}
\end{figure}

% -----------------------------------------------------------------------
\section{Başlangıç Topolojisi}
\label{sec:initial}

$X$ kümesi ve haritalar ailesi $\{f_\alpha : X \to (Y_\alpha, \tau_\alpha)\}_{\alpha \in A}$
verilsin.

\begin{definition}[Başlangıç Topolojisi]
\label{def:initial_topology}
\textbf{Başlangıç topolojisi} $\tau_{\mathrm{ini}}$, her $f_\alpha$'yı $(X, \tau_{\mathrm{ini}}) \to (Y_\alpha, \tau_\alpha)$
sürekli kılan topolojilerin kesişimidir. Alt-baz:
\[
\mathcal{S} = \{f_\alpha^{-1}(U) \mid U \in \tau_\alpha,\ \alpha \in A\},
\quad
\tau_{\mathrm{ini}} = \langle \mathcal{S} \rangle.
\]
\end{definition}

\begin{theorem}[Varlık ve Teklik]
\label{thm:initial_exists}
$\tau_{\mathrm{ini}}$ tekil ve vardır; her $f_\alpha$'yı sürekli kılan tüm topolojilerin
en kaba olanıdır.
\end{theorem}

\begin{proof}[Kanıt İskeleti]
$\mathcal{S}$'nin ürettiği topoloji $\langle\mathcal{S}\rangle$ her $f_\alpha$'yı
sürekli kılar (çünkü $f_\alpha^{-1}(U) \in \langle\mathcal{S}\rangle$ her $U \in \tau_\alpha$
için). Eğer $\tau'$ de her $f_\alpha$'yı sürekli kılıyorsa, $\mathcal{S} \subseteq \tau'$
olduğundan $\langle\mathcal{S}\rangle \subseteq \tau'$; yani $\tau_{\mathrm{ini}}$ en kaba.
\end{proof}

\begin{theorem}[Evrensel Özellik]
\label{thm:initial_universal}
$h: Z \to (X, \tau_{\mathrm{ini}})$ süreklidir ancak ve ancak $f_\alpha \circ h$ her
$\alpha$ için sürekli olduğunda.
\end{theorem}

\begin{proof}[İspat eskizi]
($\Rightarrow$) $h$ sürekli ve $f_\alpha$ sürekli ise bileşke $f_\alpha \circ h$
süreklidir. ($\Leftarrow$) $\tau_{\mathrm{ini}}$'nin bir alt-baz elemanı
$f_\alpha^{-1}(U)$ biçimindedir ve ön-görüntüsü
$h^{-1}(f_\alpha^{-1}(U)) = (f_\alpha \circ h)^{-1}(U)$, ki bu $f_\alpha \circ h$
sürekli olduğundan $Z$'de açıktır. Süreklilik alt-bazda denetlenebildiğinden $h$
süreklidir.
\end{proof}

\begin{figure}[ht]
\centering
\begin{tikzpicture}[font=\small, node distance=14mm,
    uz/.style={draw, rounded corners, fill=blue!6, inner sep=6pt, minimum width=10mm}]
  \node[uz] (Z) at (0,0) {$Z$};
  \node[uz] (X) at (3.4,0) {$(X,\tau_{\mathrm{ini}})$};
  \node[uz] (Y) at (3.4,-2.4) {$(Y_\alpha,\tau_\alpha)$};

  \draw[-{Stealth[length=2.6mm]}, thick, dashed]
    (Z) -- node[above, font=\scriptsize] {$h$ (tek, sürekli)} (X);
  \draw[-{Stealth[length=2.6mm]}, thick]
    (X) -- node[right, font=\scriptsize] {$f_\alpha$} (Y);
  \draw[-{Stealth[length=2.6mm]}, thick]
    (Z) -- node[below left, font=\scriptsize, pos=0.55] {$f_\alpha\circ h$} (Y);

  \node[font=\scriptsize, align=center] at (1.7,-3.5)
    {Evrensel özellik: $h$ sürekli $\iff$ her $f_\alpha\circ h$ sürekli.\\
     $\tau_{\mathrm{ini}}$ bu faktörizasyonu mümkün kılan \emph{tek} topolojidir.};
\end{tikzpicture}

\caption{Evrensel özellik: $h$ sürekli $\iff$ her $f_\alpha\circ h$ sürekli;
$\tau_{\mathrm{ini}}$ bu faktörizasyonu mümkün kılan tek topolojidir.}
\label{fig:ch13_evrensel}
\end{figure}

\begin{theorem}[Özel Haller]
\label{thm:initial_special}
\begin{enumerate}
  \item Altuzay topolojisi, ekleme haritası $i: A \hookrightarrow X$'in başlangıç
        topolojisidir.
  \item Çarpım topolojisi, projeksiyon haritaları $\{\pi_\alpha\}$'nın başlangıç
        topolojisidir.
\end{enumerate}
\end{theorem}

% -----------------------------------------------------------------------
\section{Algoritmalar}
\label{sec:initial_algorithms}

\subsection{Başlangıç Topolojisi İnşası}

\begin{algorithm}
\caption{Başlangıç Topolojisi (Sonlu Durum)}
\label{alg:initial_topology}
\begin{algorithmic}[1]
\Require Sonlu taşıyıcı $X$, haritalar $\{(f_\alpha, \tau_\alpha)\}$
\Ensure $\tau_{\mathrm{ini}}$ — en kaba süreklilik topolojisi
\State $\mathcal{S} \leftarrow \emptyset$
\For{her $f_\alpha$}
  \For{her $U \in \tau_\alpha$}
    \State $\mathcal{S} \leftarrow \mathcal{S} \cup \{f_\alpha^{-1}(U)\}$
    \State \textit{// $f_\alpha^{-1}(U) = \{x \in X \mid f_\alpha(x) \in U\}$}
  \EndFor
\EndFor
\State \Return $\langle \mathcal{S} \rangle$ \textit{// alt-bazdan topoloji üret}
\end{algorithmic}
\end{algorithm}

\textbf{Karmaşıklık.} Ön-görüntü hesabı $O(|\tau_\alpha| \cdot |X|)$; topoloji üretimi
$O(|\mathcal{S}| \cdot 2^{|X|})$ en kötü durumda.

\subsection{Son Topoloji İnşası}

\begin{algorithm}
\caption{Son Topoloji (Sonlu Durum)}
\label{alg:final_topology}
\begin{algorithmic}[1]
\Require Sonlu $Y$, kaynak uzaylar $\{(X_i, \sigma_i)\}$, haritalar $\{g_i : X_i \to Y\}$
\Ensure $\tau_{\mathrm{fin}}$ — en ince süreklilik topolojisi
\State $\tau_{\mathrm{fin}} \leftarrow \emptyset$
\For{her $V \subseteq Y$}
  \If{$g_i^{-1}(V) \in \sigma_i$ her $i$ için}
    \State $\tau_{\mathrm{fin}} \leftarrow \tau_{\mathrm{fin}} \cup \{V\}$
  \EndIf
\EndFor
\State \Return $\tau_{\mathrm{fin}}$
\end{algorithmic}
\end{algorithm}

\textbf{Karmaşıklık.} $O(2^{|Y|} \cdot \sum_i |\sigma_i|)$.

% -----------------------------------------------------------------------
\section{pytop API}
\label{sec:initial_api}

\begin{lstlisting}[language=Python, caption={Başlangıç Topolojisi API}]
from pytop.finite_map_analysis import FiniteMap
from pytop import (
    initial_topology_from_maps,
    generic_quotient_space_from_map,
    discrete_topology,
    indiscrete_topology,
    make_topology,
    sierpinski_space,
    finite_subspace,
    binary_product,
    is_continuous_finite_map,
    is_t0,
    is_t2,
)

# FiniteMap dogrudan veri sinifi olarak kullanilir (make_function degil)
# tau_ini = initial_topology_from_maps(carrier, [f1, f2, ...])
# quot    = generic_quotient_space_from_map(q_map)   # son topoloji, bolum ozel hali
\end{lstlisting}

\textbf{Önemli:} \texttt{pytop.quotient\_space\_from\_map} sembolik motoru kullanır;
sonlu-duyarlı sürüm \texttt{generic\_quotient\_space\_from\_map}'tir.

% -----------------------------------------------------------------------
\section{Örnekler}
\label{sec:initial_examples}

\subsection{Tek Harita ile Başlangıç Topolojisi}

\begin{lstlisting}[language=Python]
X_carrier = [0, 1, 2, 3]
Y_d = discrete_topology(0, 1)
X_ph = make_topology(X_carrier, set(), set(X_carrier))

f = FiniteMap(domain=X_ph, codomain=Y_d, name="f",
              mapping={0: 0, 1: 0, 2: 1, 3: 1})
tau_f = initial_topology_from_maps(X_carrier, [f])

print("Baslangic topolojisi:")
for u in sorted([sorted(u) for u in tau_f.topology], key=lambda x: (len(x), x)):
    print("  ", u if u else "empty")
print("Eleman sayisi:", len(tau_f.topology))
\end{lstlisting}

\begin{lstlisting}[style=output]
Baslangic topolojisi:
   empty
   [0, 1]
   [2, 3]
   [0, 1, 2, 3]
Eleman sayisi: 4
\end{lstlisting}

$f^{-1}(\{0\}) = \{0,1\}$, $f^{-1}(\{1\}) = \{2,3\}$ — alt-baz $\{\{0,1\},\{2,3\}\}$,
üretilen topoloji 4 elemanlı.

\subsection{İki Harita — Topoloji İncelir}

\begin{lstlisting}[language=Python]
S_sier = sierpinski_space()
g = FiniteMap(domain=X_ph, codomain=S_sier, name="g",
              mapping={0: 0, 1: 1, 2: 1, 3: 1})
tau_fg = initial_topology_from_maps(X_carrier, [f, g])

print("f tek basina:", len(tau_f.topology), "eleman  |  f+g:", len(tau_fg.topology), "eleman")
\end{lstlisting}

\begin{lstlisting}[style=output]
f tek basina: 4 eleman  |  f+g: 6 eleman
\end{lstlisting}

$g^{-1}(\{1\}) = \{1,2,3\}$ yeni bir alt-baz elemanı ekler; kesişim $\{0,1\} \cap \{1,2,3\} = \{1\}$
yeni bir açık küme üretir.

\subsection{Altuzay = Başlangıç}

\begin{lstlisting}[language=Python]
fts_X = make_topology([0, 1, 2, 3], set(), {0, 1}, {2, 3}, {0, 1, 2, 3})
sub_A  = finite_subspace(fts_X, [1, 2])

A_dom  = discrete_topology(1, 2)
inc    = FiniteMap(domain=A_dom, codomain=fts_X, name="i", mapping={1: 1, 2: 2})
init_A = initial_topology_from_maps([1, 2], [inc])

sub_tau  = sorted([sorted(u) for u in sub_A.topology],  key=lambda x: (len(x), x))
init_tau = sorted([sorted(u) for u in init_A.topology], key=lambda x: (len(x), x))
print("Esit mi:", sub_tau == init_tau)
\end{lstlisting}

\begin{lstlisting}[style=output]
Esit mi: True
\end{lstlisting}

\subsection{Çarpım = Başlangıç}

\begin{lstlisting}[language=Python]
d2   = discrete_topology(0, 1)
prod = binary_product(d2, d2)

pi_x = FiniteMap(domain=prod, codomain=d2, name="pi_x",
                 mapping={(0,0):0, (0,1):0, (1,0):1, (1,1):1})
pi_y = FiniteMap(domain=prod, codomain=d2, name="pi_y",
                 mapping={(0,0):0, (0,1):1, (1,0):0, (1,1):1})

init_prod = initial_topology_from_maps(list(prod.carrier), [pi_x, pi_y])
prod_fs   = {frozenset(u) for u in prod.topology}
init_fs   = {frozenset(u) for u in init_prod.topology}
print("Carpim == Baslangic:", prod_fs == init_fs)
\end{lstlisting}

\begin{lstlisting}[style=output]
Carpim == Baslangic: True
\end{lstlisting}

\subsection{Son Topoloji — Manuel Hesap}

\begin{lstlisting}[language=Python]
X1     = make_topology([0, 1, 2], set(), {0, 1}, {0, 1, 2})
Y_car  = [0, 1]
g1_map = {0: 0, 1: 0, 2: 1}
tau_X1 = {frozenset(u) for u in X1.topology}

open_Y = []
for mask in range(1 << len(Y_car)):
    V        = frozenset(Y_car[i] for i in range(len(Y_car)) if mask & (1 << i))
    preimage = frozenset(x for x in X1.carrier if g1_map[x] in V)
    if preimage in tau_X1:
        open_Y.append(set(V))

final_Y   = make_topology(Y_car, *open_Y)
final_tau = sorted([sorted(u) for u in final_Y.topology], key=lambda x: (len(x), x))
print("Son topoloji:", final_tau)
\end{lstlisting}

\begin{lstlisting}[style=output]
Son topoloji: [[], [0], [0, 1]]
\end{lstlisting}

\subsection{Bölüm = Son Topoloji}

\begin{lstlisting}[language=Python]
Y_ind    = indiscrete_topology(0, 1)
q        = FiniteMap(domain=X1, codomain=Y_ind, name="q",
                     mapping={0: 0, 1: 0, 2: 1})
quot_Y   = generic_quotient_space_from_map(q)
quot_tau = sorted([sorted(u) for u in quot_Y.topology], key=lambda x: (len(x), x))

print("Bolum topolojisi:", quot_tau)
print("Son topoloji ile esit:", quot_tau == final_tau)
\end{lstlisting}

\begin{lstlisting}[style=output]
Bolum topolojisi: [[], [0], [0, 1]]
Son topoloji ile esit: True
\end{lstlisting}

\begin{figure}[ht]
\centering
\begin{tikzpicture}[font=\small]
  % Source X_1 (top left)
  \draw[rounded corners=6pt, thick] (-0.4,0.5) rectangle (1.8,1.9);
  \node at (0.7,2.15) {$(X_1,\sigma_1)$};
  \draw[blue!70!black, thick, fill=blue!12]
    (0.7,1.2) ellipse [x radius=0.55, y radius=0.32];

  % Source X_2 (bottom left)
  \draw[rounded corners=6pt, thick] (-0.4,-1.9) rectangle (1.8,-0.5);
  \node at (0.7,-2.2) {$(X_2,\sigma_2)$};
  \draw[red!70!black, thick, fill=red!12]
    (0.7,-1.2) ellipse [x radius=0.55, y radius=0.32];

  % Target Y (right)
  \draw[rounded corners=6pt, thick] (5.2,-1.5) rectangle (7.7,1.6);
  \node at (6.45,1.85) {$Y$ (taşıyıcı)};
  \draw[violet!75!black, thick, fill=violet!10]
    (6.45,0.0) ellipse [x radius=0.85, y radius=0.55];
  \node[violet!75!black, font=\scriptsize] at (6.45,0.0) {$V$};
  \node[violet!75!black, font=\scriptsize, align=center] at (6.45,-1.05)
    {$V$ açık $\iff$\\$g_i^{-1}(V)\in\sigma_i$};

  % Maps
  \draw[blue!70!black, -{Stealth[length=2.4mm]}, thick]
    (1.4,1.2) to[bend left=12] node[above, font=\scriptsize, pos=0.5] {$g_1$} (5.55,0.45);
  \draw[red!70!black, -{Stealth[length=2.4mm]}, thick]
    (1.4,-1.2) to[bend right=12] node[below, font=\scriptsize, pos=0.5] {$g_2$} (5.55,-0.45);

  \node[font=\scriptsize, align=center] at (3.3,-2.9)
    {İleri itme: $\tau_{\mathrm{fin}}$ her $g_i$'yi sürekli kılan \emph{en ince} topoloji;\\
     bölüm topolojisi tek surjektif $q$'nun son topolojisidir.};
\end{tikzpicture}

\caption{Son topoloji: kaynak uzayları ortak bir hedefe iten en ince topoloji;
bölüm topolojisi tek surjektif haritanın son topolojisidir.}
\label{fig:ch13_son}
\end{figure}

\subsection{Yeni Örnek — ``En Kaba'' Olmanın Doğrudan Doğrulanması}

Başlangıç topolojisinin tanımı \emph{en kaba} olmaktır: $f$'yi sürekli kılar, ama
ondan daha kaba hiçbir topoloji kılamaz. Bunu süreklilik yüklemiyle doğrulayalım.

\begin{lstlisting}[language=Python]
ini      = [set(u) for u in tau_f.topology]
codom    = [set(u) for u in Y_d.topology]
coarser  = [set(), set(X_carrier)]   # indiscrete: tau_ini'den daha kaba

print("f, tau_ini'ye gore surekli mi:",
      is_continuous_finite_map(X_carrier, ini, list(Y_d.carrier), codom, f.mapping))
print("f, daha kaba (indiscrete) topolojiye gore surekli mi:",
      is_continuous_finite_map(X_carrier, coarser, list(Y_d.carrier), codom, f.mapping))
print("tau_ini eleman sayisi:", len(tau_f.topology))
\end{lstlisting}

\begin{lstlisting}[style=output]
f, tau_ini'ye gore surekli mi: True
f, daha kaba (indiscrete) topolojiye gore surekli mi: False
tau_ini eleman sayisi: 4
\end{lstlisting}

$\tau_{\mathrm{ini}}$ ile $f$ sürekli; aynı $f$ indiscrete topolojiyle sürekli
\textbf{değil} — çünkü $f^{-1}(\{0\})=\{0,1\}$ orada açık değildir. $\tau_{\mathrm{ini}}$
tam da bu açık kümeyi ekleyecek kadar incedir, fazlası değil.

\begin{karsiornek}
``Daha çok açık küme her zaman daha iyidir'' sanmayın. Discrete topoloji de $f$'yi
sürekli kılar (16 açık küme), ama \emph{en kaba} değildir — başlangıç topolojisi
yalnızca süreklilik için gereken 4 açık kümeyi koyar. Gereğinden ince bir topoloji
evrensel özelliği bozar: $\tau_{\mathrm{ini}}$'den daha ince bir $\tau'$ seçerseniz,
$f_\alpha\circ h$ sürekli olduğu hâlde $h$ sürekli olmayan bir $h$ bulunabilir.
\end{karsiornek}

\subsection{Yeni Örnek — Evrensel Özelliğin Doğrulanması}

İki harita $f,g: X \to \mathrm{discrete}(\{0,1\})$ ile $X=\{a,b\}$ üzerinde başlangıç
topolojisini kurar, ardından $h: Z \to X$ sürekliliğinin her bileşenin sürekliliğine
denk olduğunu gösteririz.

\begin{lstlisting}[language=Python]
Xc   = ["a", "b"]
Xph  = make_topology(Xc, set(), set(Xc))
Yd2  = discrete_topology(0, 1)
fA   = FiniteMap(domain=Xph, codomain=Yd2, name="fA", mapping={"a": 0, "b": 1})
gA   = FiniteMap(domain=Xph, codomain=Yd2, name="gA", mapping={"a": 1, "b": 0})

tau_X  = initial_topology_from_maps(Xc, [fA, gA])
Xopens = [set(u) for u in tau_X.topology]
Yopens = [set(u) for u in Yd2.topology]

Zc, Zopens = ["p", "q"], [set(), {"p"}, {"q"}, {"p", "q"}]
h  = {"p": "a", "q": "b"}
fh = {z: fA.mapping[h[z]] for z in Zc}     # fA . h
gh = {z: gA.mapping[h[z]] for z in Zc}     # gA . h

c_fh = is_continuous_finite_map(Zc, Zopens, list(Yd2.carrier), Yopens, fh)
c_gh = is_continuous_finite_map(Zc, Zopens, list(Yd2.carrier), Yopens, gh)
c_h  = is_continuous_finite_map(Zc, Zopens, Xc, Xopens, h)

print("tau_ini(X) eleman sayisi:", len(tau_X.topology))
print("fA.h surekli:", c_fh, "| gA.h surekli:", c_gh, "| h surekli:", c_h)
print("Evrensel ozellik saglaniyor mu:", c_h == (c_fh and c_gh))
\end{lstlisting}

\begin{lstlisting}[style=output]
tau_ini(X) eleman sayisi: 4
fA.h surekli: True | gA.h surekli: True | h surekli: True
Evrensel ozellik saglaniyor mu: True
\end{lstlisting}

$\tau_{\mathrm{ini}}(X)$ discrete$(\{a,b\})$'ye eşittir; $h$ süreklidir ancak ve ancak
hem $f_A\circ h$ hem $g_A\circ h$ sürekli olduğunda. Yüklem bu denkliği
\texttt{c\_h == (c\_fh and c\_gh)} ile doğrular.

% -----------------------------------------------------------------------
\section{Alıştırmalar}
\label{sec:initial_exercises}

\begin{enumerate}
  \item $f: \{a,b,c,d\} \to \mathrm{discrete}(\{0,1\})$, $f(a)=f(b)=0$, $f(c)=f(d)=1$
        için \texttt{initial\_topology\_from\_maps} kullanın ve sonucu
        \texttt{finite\_subspace} ile karşılaştırın.

  \item Çarpım özel halinde yalnızca $\pi_x$ kullanıldığında başlangıç topolojisi
        ne olur? Discrete'den daha kaba mı?

  \item $X=\{0,1,2,3,4\}$, $\sigma_X = $ discrete, $g: X \to \{0,1,2\}$ haritası
        için son topolojiyi manuel hesaplayın.

  \item \textit{(Teori)} $h: Z \to (X, \tau_{\mathrm{ini}})$ sürekliliğinin
        $f_\alpha \circ h$ sürekliliğine denk olduğunu ispatlayın.

  \item \textit{(Teori)} $\tau_{\mathrm{fin}}$'den daha ince hiçbir topoloji her
        $g_i$'yi sürekli kılamaz; bunu tanımdaki en-ince koşulundan türetin.

  \item \texttt{initial\_topology\_from\_maps} ile bir
        $f: \{0,1,2,3\} \to \mathrm{discrete}(\{0,1\})$, $f(\{0,1\})=0$,
        $f(\{2,3\})=1$ haritası için $\tau_{\mathrm{ini}}$ kurun. Sonra
        \texttt{is\_continuous\_finite\_map} ile $f$'nin $\tau_{\mathrm{ini}}$'ye göre
        sürekli ama indiscrete topolojiye göre sürekli olmadığını doğrulayın. Aynı
        $f$ discrete topolojiyle de sürekli midir? Bu, ``en kaba'' iddiasını nasıl
        destekler?

  \item \textit{(Evrensel özellik)} $X=\{a,b\}$ üzerinde
        $f(a)=0,f(b)=1$ ve $g(a)=1,g(b)=0$ (kodomen \texttt{discrete(\{0,1\})}) için
        $\tau_{\mathrm{ini}}$'yi hesaplayın. $Z=\mathrm{discrete}(\{p,q\})$,
        $h(p)=a$, $h(q)=b$ alın. \texttt{is\_continuous\_finite\_map} ile
        $f\circ h$, $g\circ h$ ve $h$ sürekliliklerini hesaplayıp $h$'nin
        sürekliliğinin bileşenlerin sürekliliğine denk olduğunu
        (\texttt{c\_h == (c\_fh and c\_gh)}) gösterin.
\end{enumerate}
